%% DetailedBalanceChecker — expanded manuscript draft
%% Target: J. Chem. Theory Comput.
%% Author: S. Whitby
%%
%% Build: pdflatex paper && bibtex paper && pdflatex paper && pdflatex paper
%% (or use latexmk -pdf paper)

\documentclass[twocolumn,10pt]{article}
\usepackage[top=1.8cm,bottom=2.0cm,left=1.4cm,right=1.4cm]{geometry}
\usepackage{amsmath,amssymb}
\usepackage{graphicx}
\usepackage{algorithm}
\usepackage{algpseudocode}
\usepackage{booktabs}
\usepackage{microtype}
\usepackage{xcolor}
\usepackage[colorlinks=true,linkcolor=blue,citecolor=blue,urlcolor=blue]{hyperref}
\usepackage{caption}
\usepackage{listings}

\lstset{
  basicstyle=\ttfamily\footnotesize,
  breaklines=true,
  frame=single,
  framesep=3pt,
  xleftmargin=5pt,
  xrightmargin=5pt
}

\setlength{\columnsep}{0.55cm}
\captionsetup{font=small, labelfont=bf}

% Lattice-site box
\newcommand{\lsite}[1]{\fbox{\ensuremath{\,#1\,}}}
% Metropolis min-expression shorthand
\newcommand{\Mf}[1]{\min\!\bigl(1,e^{#1}\bigr)}
% Tool name shorthand
\newcommand{\DBC}{\textit{DetailedBalanceChecker}}

\title{\textbf{Symbolic Verification of Detailed Balance\\[-2pt] in Monte Carlo Algorithms}}
\author{S.~Whitby}
\date{}

\begin{document}
\maketitle
\thispagestyle{empty}

% ============================================================
\begin{abstract}
\noindent
We present \DBC{}, a computer-algebra framework for the exact, symbolic
verification of detailed balance in Monte Carlo (MC) algorithms.  Every call
to the random-number generator is intercepted and mapped to reads from a
binary tape, enabling a breadth-first search (BFS) over all execution paths
of the algorithm.  The resulting symbolic transition matrix is tested against
the detailed-balance condition algebraically, treating coupling constants,
field functions, and inverse temperature as free symbols.  A PASS constitutes
a mathematical proof valid for all parameter values simultaneously; a FAIL
identifies the exact symbolic expressions that violate balance for each
offending state pair.  The primary verification strategy exploits the
algebraic structure of Metropolis-type acceptance: after expanding over sign
cases of the coupling constants, each detailed-balance expression reduces to
a sum of terms with common exponential factors, verified in microseconds by
polynomial coefficient comparison.  A canonical-ordering oracle and
generator-based symmetry check enable exploitation of the lattice symmetry
group $G$ to compress the BFS and the verification, yielding near-$|G|$-fold
reductions in cost.  Ergodicity is checked independently via BFS state-count
comparison.  Physical fidelity metrics quantify how closely the MC transition
matrix matches reference Glauber dynamics.  We demonstrate the framework on
1D Kawasaki dynamics, a rightward-only cluster algorithm that passes every
numerical statistical test while violating detailed balance, and 2D Virtual
Move Monte Carlo (VMMC) in both correct and broken variants.  Intended use
cases include regression testing, algorithm development, and as a formal
oracle in an LLM-driven algorithm discovery loop.  The implementation is
available at \url{https://github.com/Sam-Whitby/CheckDetailedBalance2}.
\end{abstract}

% ============================================================
\section{Introduction}

Monte Carlo (MC) simulation is the primary tool for sampling equilibrium
distributions in statistical physics and soft
matter~\cite{Metropolis1953,Frenkel2002,Allen1987}.  The correctness of a MC
algorithm rests on two conditions: ergodicity (every state is reachable from
every other) and the \emph{detailed balance} (DB) condition
\begin{equation}
  T(i \to j)\,\pi(i) \;=\; T(j \to i)\,\pi(j),
  \label{eq:db}
\end{equation}
where $T(i\to j)$ is the transition rate from microstate $i$ to $j$ and
$\pi(i)\propto e^{-\beta E(i)}$ is the target Boltzmann weight.  Together
with ergodicity, Eq.~\eqref{eq:db} guarantees convergence to the correct
equilibrium distribution~\cite{Levin2009}.

Algorithm correctness is overwhelmingly assessed by numerical tests.  The
most common approach runs a long Markov chain and compares the empirical
stationary distribution to the Boltzmann distribution via the
Kullback--Leibler (KL) divergence~\cite{Newman1999} or a $\chi^2$
goodness-of-fit statistic.  The joint-distribution test of
Geweke~\cite{Geweke2004} is the most principled numerical approach,
comparing forward and time-reversed marginals of successive chain states.
These tests share a fundamental limitation: they are statistical and
therefore cannot certify correctness.  They can only fail to detect a
violation.  An algorithm whose stationary distribution happens to be
Boltzmann-distributed under the specific energy function and parameter values
used in the test will pass every numerical test, even if it violates detailed
balance and would sample incorrectly for a different energy~--- see
Section~\ref{sec:cluster} for a concrete example.

Analytical proofs of detailed balance for specific algorithms---Metropolis--Hastings~\cite{Hastings1970},
VMMC~\cite{Whitelam2007,Whitelam2009}, Swendsen--Wang~\cite{SwendsenWang1987},
Wolff~\cite{Wolff1989}---are algorithm-specific and laborious.  They are
also susceptible to the same sign errors and missed edge cases that numerical
tests at benign parameters cannot detect.  For complex cluster algorithms,
where the acceptance criterion emerges from a multi-step cluster-growth
procedure, the analytical proof requires careful book-keeping of all path
weights, and subtle implementation errors may not be caught until the
algorithm is applied to a model where the violation is not hidden by energy
symmetries.

A more general reversibility condition---the Kolmogorov cycle
criterion~\cite{Kelly1979}---states that the chain is reversible if and only
if the product of transition rates is equal in both directions around every
closed loop of states.  Checking this without knowing the stationary
distribution requires either enumerating all cycles or matrix
diagonalisation at $\mathcal{O}(N^3)$ cost~\cite{Norris1998}, which becomes
algebraically intractable for symbolic parameter values.  When the energy $E$
is known, as it always is in physics simulations, detailed balance reduces
instead to $N(N-1)$ pairwise checks of Eq.~\eqref{eq:db}---a tractable
algebraic problem.

We present an automated framework for the \emph{exact symbolic} verification
of Eq.~\eqref{eq:db}.  The framework intercepts every random call in the
algorithm, constructs an exact symbolic transition matrix by BFS over all
execution paths, and verifies detailed balance algebraically.  All coupling
constants and field functions are treated as free real symbols, so a PASS
constitutes a mathematical proof valid for all parameter values
simultaneously, including parameter combinations not accessible to numerical
testing.  A FAIL yields exact symbolic expressions identifying which
transitions violate detailed balance and by how much, enabling diagnosis and
correction.  We also describe how symmetry of the state space under a lattice
symmetry group $G$ can be exploited to achieve near-$|G|$-fold speedups in
both the BFS and the verification step, extending the practical reach of the
method.

% ============================================================
\section{Related Work}
\label{sec:related}

The dominant paradigm for MC correctness checking remains numerical.
Histogram comparison against the Boltzmann distribution via KL divergence or
$\chi^2$ statistics~\cite{Newman1999} is standard practice, and the
joint-distribution test of Geweke~\cite{Geweke2004}---which compares forward
and time-reversed marginals of successive chain states---is the most
principled purely numerical approach.  Both are statistical.

Formal methods for Markov chains have advanced considerably through
\emph{parametric model checking}~\cite{Daws2004}, implemented in PRISM~\cite{Kwiatkowska2011}
and Storm~\cite{Dehnert2017}.  In principle, steady-state probabilities
computed by these tools could be combined with transition-matrix entries to
check detailed balance.  However, PRISM and Storm operate on models specified
in dedicated input languages.  Translating a Monte Carlo algorithm into such
a format requires essentially the same case analysis as writing an analytic
proof, and must be redone for each algorithmic variant.  Detailed balance is
not a native property of any standard temporal logic, so any verification
requires additional post-processing outside the model-checker.

Exact symbolic inference for probabilistic programs---PSI~\cite{Gehr2016}
and similar tools---can in principle yield a single transition probability
$T(i\to j)$ for a specific state pair.  However, these systems target
Bayesian inference, not exhaustive relational checking across all state
pairs.

It is worth noting that detailed balance is \emph{sufficient} but not
\emph{necessary} for a Boltzmann stationary distribution.  The weaker global
balance condition suffices~\cite{Manousiouthakis1999}, and algorithms such as
event-chain Monte Carlo (ECMC)~\cite{Bernard2009}---which achieves
rejection-free updates with correct equilibrium statistics by satisfying
only global balance via lifting variables---would \emph{correctly} fail the
detailed-balance checker.  This is a feature: it defines a precise and
demanding search objective.

To the authors' knowledge, no existing framework provides automated
detailed-balance checking as a first-class operation for Monte Carlo
algorithms written as executable code, nor verifies correctness across the
full parameter space in a single algebraic pass.

% ============================================================
\section{Theoretical Framework}
\label{sec:method}

We describe the algorithm in abstract terms, independent of any particular
implementation.

\subsection{Problem formulation}

We consider a finite set of microstates $\mathcal{S}$ with a user-supplied
energy function $E:\mathcal{S}\to\mathbb{R}$.  An MC algorithm is a
function $\mathcal{A}:\mathcal{S}\times\{0,1\}^\ast\to\mathcal{S}$ mapping
an initial state and a binary tape to a final state, where each random call
inside $\mathcal{A}$ reads one or more bits from the tape.  The transition
probability from state $i$ to state $j$ is
\begin{equation}
  T(i\to j) = \sum_{\mathbf{b}:\,\mathcal{A}(i,\mathbf{b})=j} w(\mathbf{b}),
\end{equation}
where $w(\mathbf{b})$ is the probability of the bit sequence $\mathbf{b}$.
The goal is to verify Eq.~\eqref{eq:db} symbolically, treating all model
parameters (coupling constants, field values, inverse temperature $\beta$) as
free real symbols.

\subsection{State enumeration via \textsc{BitsToState}}
\label{sec:bitstostate}

The framework requires a user-supplied injective map
$\textsc{BitsToState}:\{0,1\}^\ast\to\mathcal{S}\cup\{\varnothing\}$ from
binary strings to states, returning $\varnothing$ for strings that do not
correspond to valid states.  This induces a total ordering of states by
bit-string length: shorter strings enumerate simpler states first.  The
checker iterates over all binary strings up to a user-specified maximum
length, calling $\textsc{BitsToState}$ on each, and launches a BFS
from each valid seed state.  Any newly discovered state becomes an
additional seed.

For 2D square-lattice systems, the reference implementation provides a
bijective encoding of labeled-particle configurations into integers: a
non-negative integer $\mathrm{ID}$ encodes (i) the lattice size $L^2$, via
a prefix count over all arrays of smaller size; (ii) the particle number
$N$, via a secondary prefix; (iii) the sorted particle positions, ranked by
the combinatorial number system; and (iv) the assignment of particle types to
positions, ranked by the factorial number system.  Each step has an explicit
inverse, so any integer can be decoded to the corresponding state.
$\textsc{BitsToState}$ interprets a bit string as the binary representation
of $\mathrm{ID}$ and calls the decoder.  Additional validity filters (e.g.,
requiring $L^2$ to be a perfect square for 2D grids) return $\varnothing$
for invalid IDs, causing the checker to skip those strings at zero BFS cost.

\subsection{Reducing randomness to a binary tape}
\label{sec:bits}

Every random call inside the algorithm is intercepted and replaced by a
deterministic token that reads from a global binary tape
$\mathbf{b}=b_1b_2\cdots\in\{0,1\}^\ast$.  Four interception rules cover all
supported call types.

\smallskip\noindent\textbf{Discrete uniform integers.}
Sampling from $\{l,\ldots,h\}$ uses rejection sampling over blocks of
$c=\lceil\log_2(h-l+1)\rceil$ bits.  A block $B$ maps to outcome $l+B$ if
$B < h-l+1$; otherwise a new block is read.  Summing over all retry paths
gives each outcome weight $1/(h-l+1)$ exactly, matching the target uniform
distribution.  This handles \texttt{RandomInteger[\{lo,hi\}]} and uniform
\texttt{RandomChoice} over a list.

\smallskip\noindent\textbf{Continuous uniform reals (interval tracking).}
A call $U\sim\mathrm{Uniform}[0,1]$ creates a symbolic token representing a
latent variable $U$ on the interval $[0,1]$.  Each subsequent comparison $U
< p$ reads one bit: bit~0 restricts $U$ to $[0,p]$ with probability weight
$p/(h-l)$ (where $[l,h]$ is the current interval); bit~1 restricts to
$[p,h]$ with weight $(h-p)/(h-l)$.  Successive comparisons on the same
variable---possibly to different thresholds in different branches---are all
conditioned correctly on all prior constraints, implementing exact
inverse-CDF sampling.  This handles \texttt{RandomReal[]} and is critical
for Metropolis acceptance steps, where a single $U$ is compared to a single
threshold but that threshold may itself be symbolic.

\smallskip\noindent\textbf{Weighted discrete choice.}
Sampling from a weighted distribution
$\{(w_k, x_k)\}_{k=1}^n$ with $\sum_k w_k = W$ is decomposed into a
sequential chain of Bernoulli trials using the binary representation of the
cumulative distribution.  Specifically, starting from the full weight sum
$W$, the first trial asks whether $U < w_1/W$: bit~0 selects $x_1$ (with
weight $w_1/W$); bit~1 continues to a trial on the remaining distribution.
Each fresh trial uses an independent interval-tracking token.  The weight of
any complete path to outcome $x_k$ is exactly $w_k/W$.  This handles
\texttt{RandomChoice[weights $\to$ elements]}.

\smallskip\noindent\textbf{Discretised Gaussian jumps.}
For algorithms on lattices where displacements are necessarily integer-valued,
a normally distributed jump of standard deviation $\sigma$ can be represented
exactly.  The probability that a continuous $\mathcal{N}(0,\sigma^2)$ random
variable falls in the unit bin centred on integer $d$ is
\begin{equation}
  w(d) = \tfrac{1}{2}\!\left[\operatorname{erf}\!\left(\tfrac{s_d+\frac{1}{2}}{\sqrt{2}\,\sigma}\right)
  - \operatorname{erf}\!\left(\tfrac{s_d-\frac{1}{2}}{\sqrt{2}\,\sigma}\right)\right],
  \label{eq:normalbin}
\end{equation}
where $s_d$ is the signed displacement on the ring.  The weights are
symmetric ($w(d)=w(-d)$), normalised, and sum to one over all integer
displacements within a truncation range.  The displacement is drawn via
weighted discrete choice with these weights, which the checker decomposes
into sequential Bernoulli trials.  The resulting symbolic transition matrix
entries contain \texttt{erf} functions with $\sigma$ as a free parameter, and
detailed balance is verified for all values of $\sigma$ simultaneously.

All four reductions are exact: integrating the branch weights over all bit
sequences recovers the original probability measure of each random call.
Calls of each type may be combined arbitrarily within a single algorithm
execution; the symbolic weight of any complete path is the product of the
weights of all individual bits read.

\subsection{BFS construction of the transition matrix}
\label{sec:bfs}

\begin{algorithm}[t]
\caption{Symbolic transition-matrix construction (BFS)}
\label{alg:bfs}
\begin{algorithmic}[1]
\Require Algorithm $\mathcal{A}$, seed state $s_0$, energy $E$
\State $\mathrm{queue}\leftarrow\{(s_0,\varepsilon,1)\}$;\;
       $T\leftarrow\mathbf{0}$;\;
       $\mathrm{seen}\leftarrow\{s_0\}$
\While{queue not empty}
  \State $(s,\mathbf{b},w)\leftarrow\mathrm{queue.pop}()$
  \State Run $\mathcal{A}(s)$ with tape $\mathbf{b}$, intercepting random calls
  \If{execution completes with output $s'$}
    \State $T[s\to s']\mathrel{+}=w$
    \If{$s'\notin\mathrm{seen}$}
      \State Add $s'$ to $\mathrm{seen}$;\; seed new BFS from $s'$
    \EndIf
  \Else{ (next bit $b_i$ needed)}
    \State Push $(s,\mathbf{b}{\cdot}0,w_0)$ and $(s,\mathbf{b}{\cdot}1,w_1)$
  \EndIf
\EndWhile
\State \Return $T$
\end{algorithmic}
\end{algorithm}

Algorithm~\ref{alg:bfs} constructs the symbolic transition matrix by
breadth-first search.  Starting from a seed state $s_0$, the algorithm
$\mathcal{A}$ is re-executed for each distinct binary prefix in the queue.
When a random call is encountered, execution branches into two continuations
(one per bit value), each pushed with its respective probability weight.
When execution completes, the output state $s'$ and accumulated weight $w$
contribute $T[s_0\to s']\mathrel{+}=w$.  Each newly discovered state
becomes an additional seed.  The procedure simultaneously constructs the
transition matrix and discovers the full connected component reachable from
any seed.

The BFS tree for a single state has depth equal to the maximum number of
random bits consumed by any complete execution path.  For a Metropolis
algorithm choosing among $n$ discrete options with Gaussian-displacement
acceptance, this is typically $\lceil\log_2 n\rceil + O(1)$ bits per
cluster step, giving a tree of size $\mathcal{O}(2^{\lceil\log_2 n\rceil})$
per state.  The total cost grows as $\mathcal{O}(N^2\cdot 2^d)$ where $N$
is the component size and $d$ is the maximum bit depth, so the method is
practical only for small systems (typically $N\lesssim 1000$, $d\lesssim 20$).

\subsection{Detailed balance verification}
\label{sec:dbcheck}

Given the symbolic transition matrix $T$ and energy function $E$, the checker
evaluates
\begin{equation}
  \Delta_{ij} = T(i{\to}j)\,e^{-\beta E(i)} - T(j{\to}i)\,e^{-\beta E(j)}
  \stackrel{?}{=} 0
  \label{eq:check}
\end{equation}
for every ordered pair $(i,j)$.  Here $\beta$, all coupling functions
$J(\cdot,\cdot,\cdot)$, field functions $f(\cdot,\cdot)$, and any other model
parameters are free real symbols.  Each distinct call $J(a,b,d^2)$ or
$f(x,y,L)$ is treated as an independent free real atom, so a PASS proves
correctness for \emph{all} possible energy models, not merely a specific
parametric form.

Three efficiency steps are applied before any algebra.

\medskip\noindent\textbf{Trivial-zero filter.}
Pairs where both $T(i\to j)$ and $T(j\to i)$ are structurally zero (never
connected by the algorithm) are skipped without any algebraic work.

\medskip\noindent\textbf{Syntactic deduplication.}
All remaining $\Delta_{ij}$ expressions are hashed before simplification.
Each unique expression is simplified exactly once; structurally identical
copies share the result.  For periodic lattices with translational symmetry,
this alone reduces thousands of unique pairs to hundreds---e.g.,
approximately $9\times$ for a $3\times 3$ torus.  With the additional
canonical-ordering oracle (Section~\ref{sec:symmetry}), G-orbit pairs
produce syntactically identical expressions, enabling a further
$\sim\!|G|$-fold collapse.

\medskip\noindent\textbf{FastChecker: polynomial coefficient test.}
The primary simplification strategy exploits the algebraic structure of
Metropolis acceptance.  After expanding a $\Delta_{ij}$ expression over sign
cases of the coupling constants via \texttt{PiecewiseExpand}, each
case reduces (within each feasible sign regime) to a sum of terms
$c_k\,e^{-\beta L_k}$, where $c_k$ is a rational coefficient and $L_k$ is a
polynomial in the coupling constants.  Detailed balance holds if and only if,
for each distinct exponent polynomial $L_k$, the sum of its coefficients
vanishes:
\begin{equation}
  \sum_{k:\,L_k = L} c_k = 0 \quad \forall L.
  \label{eq:fast}
\end{equation}
This is verified by pure polynomial arithmetic in microseconds.

Before applying the test, each Piecewise case must be checked for
\emph{feasibility}---whether there exist real values of the coupling constants
satisfying the case conditions.  The infeasible cases cannot contribute
violations.  The feasibility check uses a two-path strategy: (1) a fast
$\mathcal{O}(n)$ pattern-match identifies conditions consisting purely of
\texttt{Unequal} atoms between distinct expressions, which are always
satisfiable over the reals (each constraint is a co-dimension-1 surface;
any finite conjunction of such surfaces has full Lebesgue measure);
(2) for conditions not covered by the fast path, \texttt{FindInstance}
locates a concrete example, which is typically $10$--$100\times$ faster than
\texttt{Reduce} for polynomial-inequality conditions.

\medskip\noindent\textbf{Implicit equality extraction (deepCheck).}
For the rare case where the fast path is inconclusive after simplification,
a second pass (\texttt{deepCheck=True}) is applied.
\texttt{LogicalExpand} first expands any negated conjunctions in the
Piecewise conditions (e.g., $\neg(A\wedge B)\to\neg A\vee\neg B$), exposing
implicit equalities hidden inside negated conditions.
\texttt{SubstEqualities} then substitutes these equalities into the residual
expression; algebraic cancellation then confirms zero without further calls
to \texttt{FullSimplify}.  This step was required to avoid a false positive
in the VMMC algorithm at $3\times3$ grid size: without it, \texttt{FullSimplify}
treated each Piecewise case independently and missed the implicit constraint
$J_{231} = J_{232}$ implied by prior-case exclusion, erroneously reporting
1296 spurious violations.

\medskip\noindent\textbf{FullSimplify fallback.}
Any expression for which the polynomial test and \texttt{deepCheck} are both
inconclusive falls back to
$\texttt{FullSimplify}[\texttt{PiecewiseExpand}[\Delta_{ij}], \{\beta>0,\ldots\}]$
with sign assumptions on coupling constants.  This is always exact, at the
cost of higher computational time (typically $0.3$--$5$\,s per expression).

\subsection{Symmetry-aware checking}
\label{sec:symmetry}

For algorithms on periodic lattices whose energy and proposal are invariant
under a lattice symmetry group $G$, the detailed-balance expressions
$\Delta_{ij}$ and $\Delta_{i'j'}$ are equal whenever $(i',j')$ is a
$G$-orbit image of $(i,j)$.  Checking one representative per orbit would
suffice, giving an $|G|$-fold speedup.  For the square torus with full D4
symmetry (8 rotations and reflections) combined with $(L\times L)$
translations, $|G|=8L^2$, making the speedup significant for any lattice
of practical size.

The challenge is that $G$-orbit pairs reference different site indices and
different neighbour-list orderings, producing syntactically distinct symbolic
expressions even when mathematically equal.  The existing syntactic
deduplication therefore does not automatically group them.  We solve this
by making the BFS itself produce \emph{syntactically identical} expressions
for $G$-orbit pairs.

\medskip\noindent\textbf{Canonical-ordering oracle.}
The seqBernoulli decomposition tree produced by the cluster-building routine
is syntactically determined by the \emph{order} in which candidate neighbour
sites are visited.  For $G$-invariant algorithms, the topological key
\[
  (d^2(p,q),\; d^2(p_\mathrm{post},q),\; d^2(p_\mathrm{rev},q),\; \text{type}(q))
\]
where $d^2$ denotes the minimum-image squared distance, is preserved by every
$G$-isometry.  Sorting candidates by this key assigns the same canonical
order to all $G$-orbit-related states.  The BFS then produces the same
symbolic tree for every state in a $G$-orbit, making the resulting
$\Delta_{ij}$ expressions syntactically identical and hash-equivalent.
The existing syntactic deduplication then collapses the entire orbit to one
representative evaluated once, at no additional verification cost.

The canonical oracle is exposed through the standard interface function
$\texttt{\$vmmcCandidates}$ defined in the shared geometry layer.  The
checker substitutes the canonical implementation via Mathematica's
\texttt{Block} during BFS: the algorithm requires no modification and
remains unaware of the substitution:
\begin{lstlisting}
$bfsAlg = Function[s,
  Block[{$vmmcCandidates = $dbcCanonicalCandidates},
    $alg[s]]]
\end{lstlisting}

\medskip\noindent\textbf{G-invariance verification.}
After BFS, a separate check verifies that the computed $T$ matrix is indeed
$G$-invariant: $T(s\to s') = T(g(s)\to g(s'))$ for all generators $g$ and
all state pairs.  The group $G = (\mathbb{Z}/L\mathbb{Z})^2\rtimes D_4$ on an
$L\times L$ torus is generated by four elements: translate$(1,0)$,
translate$(0,1)$, rotate$_{90}$, and reflect (main diagonal).  G-invariance
under any product of generators follows by induction if it holds for all four;
thus four sweeps over all state pairs suffice to certify invariance under all
$|G|\leq 8L^2$ elements.  Each generator's action on all states is
precomputed once into a hash table, making per-pair lookups $\mathcal{O}(1)$.
The comparison is syntactic equality (valid because the canonical ordering
makes $G$-related entries syntactically identical).

If the $G$-invariance check fails---indicating the declared symmetry group
is incorrect for this algorithm---the checker reports FAIL and the DB result
is unaffected (it is always correct regardless of the symmetry declaration).
If it passes, the user has both a formal proof of DB and a certificate of
$G$-invariance.

Table~\ref{tab:symmetry} summarises the safety guarantees.

\begin{table}[t]
  \centering
  \caption{Safety guarantees for symmetry-aware checking.  The DB check
    result is always correct.  G-invariance failure warns that the declared
    symmetry is wrong; it does not produce false DB results.}
  \label{tab:symmetry}
  \small
  \setlength{\tabcolsep}{3pt}
  \begin{tabular}{@{}p{3.0cm}cc@{}}
    \toprule
    Scenario & DB & G-inv \\
    \midrule
    Satisfies DB and $G$ & PASS & PASS \\
    Violates DB & FAIL & PASS/FAIL \\
    Satisfies DB but not $G$ & PASS & FAIL \\
    Wrong $\$\mathrm{symmetryGroup}$ declaration & PASS & FAIL \\
    \bottomrule
  \end{tabular}
\end{table}

\medskip\noindent\textbf{Conditions for a valid symmetry declaration.}
Three structural conditions must hold for the $G$-invariance declaration to
be correct and the orbit reduction to give a valid speedup:
(1) seed selection is uniform over occupied sites with no position-dependent
weighting;
(2) the displacement list contains every $(dx,dy)$ paired with $(-dx,-dy)$
(proposal symmetry $P(\mathbf{d})=P(-\mathbf{d})$);
(3) the energy uses only minimum-image periodic distances (no external field).
If any condition fails, the $G$-invariance check will detect and report it.

\subsection{Ergodicity checking}
\label{sec:ergodicity}

Detailed balance alone does not guarantee a correct sampler; the chain must
also be ergodic (every state reachable from every other state).  The checker
verifies ergodicity independently of detailed balance by comparing the number
of states discovered in each BFS-connected component with the theoretical
state count.  For $N$ labeled particles on $S$ sites (with at most one
particle per site), the number of reachable configurations from any starting
configuration under a swap-type move is the falling factorial $(S)_N =
S(S-1)\cdots(S-N+1)$.  A component whose BFS-discovered state count equals
$(S)_N$ is ergodic; a smaller count signals non-ergodicity.

This check is implemented by the \texttt{ValidStateIDs} function in each
algorithm file, which specifies the theoretical count for the given system
size, and by comparing BFS-discovered state counts to this value.  A separate
ergodicity check based on the support matrix $(I+S)^{n-1}$---which has a
strictly positive $(i,j)$ entry if and only if state $j$ is reachable from
$i$~\cite{HornJohnson1985}---is available for small components where
matrix-power computation is feasible.

Notably, detailed balance does not imply ergodicity and vice versa.  A
non-ergodic algorithm can satisfy DB within each disconnected component while
failing to reach all of state space; this is demonstrated explicitly in
Section~\ref{sec:examples}.

\subsection{Physical fidelity metrics}
\label{sec:fidelity}

Beyond the binary DB and ergodicity checks, the checker computes two optional
scalar metrics measuring how closely the algorithm's transition matrix $T_\mathrm{MC}$
matches a physical reference dynamics.  For a 2D periodic square lattice, the
physical reference is Glauber single-particle dynamics~\cite{Glauber1963}.
For each state pair $(i,j)$ differing by exactly one nearest-neighbour hop
of a single particle:
\begin{equation}
  T_\mathrm{phys}(i\to j) = \frac{1}{N} \cdot \frac{1}{z}
  \cdot \frac{1}{1+e^{\beta\Delta E}},
  \label{eq:glauber}
\end{equation}
where $z=4$ is the 2D coordination number and $\Delta E = E(j)-E(i)$.
All other entries of $T_\mathrm{phys}$ are zero (diagonal entries enforce
row-sum normalisation).  $T_\mathrm{phys}$ is constructed exactly from the
enumerated state space without simulation.

\medskip\noindent\textbf{M1: relaxation spectrum ratio.}
\begin{equation}
  M_1 = \left|\frac{\lambda_2^\mathrm{MC}}{\lambda_3^\mathrm{MC}} -
  \frac{\lambda_2^\mathrm{phys}}{\lambda_3^\mathrm{phys}}\right|
  \label{eq:M1}
\end{equation}
where eigenvalues are sorted by real part (descending), $\lambda_1=1$ for any
ergodic chain.  $M_1$ measures whether the algorithm reproduces the correct
hierarchy of relaxation timescales.  $M_1=0$ indicates the ratio of the two
slowest relaxation rates matches Glauber exactly; larger $M_1$ indicates a
distorted timescale structure.

\medskip\noindent\textbf{M2: row-normalised KL divergence (fidelity score).}
\begin{equation}
  M_2 = -\sum_i \pi_i \sum_j T_\mathrm{phys}(i\to j)\,
  \log\frac{T_\mathrm{phys}(i\to j)}{T_\mathrm{MC}(i\to j)},
  \label{eq:M2}
\end{equation}
the $\pi$-weighted average KL divergence of the physical transition
distribution from the MC distribution, with the sign convention that $M_2=0$
is a perfect match and $M_2<0$ indicates departure.  A value of $-\infty$
(hard failure) indicates the algorithm cannot execute a transition that is
physically required, $T_\mathrm{MC}(i\to j)=0$ for some $(i,j)$ with
$T_\mathrm{phys}(i\to j)>0$.

These metrics are not pass/fail criteria.  A more negative $M_2$ for a
cluster algorithm such as VMMC is expected and reflects algorithmic
acceleration: the algorithm decorrelates faster than physical Brownian
dynamics by construction.  The metrics provide a quantitative basis for
comparing the \emph{kinetic character} of different MC moves targeting the
same equilibrium distribution.

\subsection{Limitations}
\label{sec:limitations}

\medskip\noindent\textbf{Termination.}
There is no guarantee that the BFS terminates for an arbitrary algorithm.
Algorithms generating unbounded random sequences correspond to infinite binary
trees; no automated procedure can decide termination in general (the halting
problem~\cite{Turing1936}).  BFS branches are truncated at a user-specified
maximum depth (default 20 bits); paths reaching this depth are silently
excluded.  A warning is printed whenever truncation occurs; false negatives
(missed violations) are possible in this regime.

\medskip\noindent\textbf{Energy function.}
The checker uses \texttt{energy[state]} exactly as written.  An incorrect
energy function invalidates both the symbolic and numerical checks
simultaneously.  Random calls inside the energy function are detected and
flagged, but other incorrectness (wrong formula, wrong sign) is invisible to
the framework.

\medskip\noindent\textbf{Hidden state.}
If the algorithm uses mutable global state, system time, or external calls
that are not intercepted, the recorded execution paths may not match actual
execution.  \texttt{CheckAlgorithmSafety} detects known unsafe call patterns
but cannot detect custom C extensions or arbitrary I/O.

Table~\ref{tab:failmodes} summarises the known false-negative and
false-positive risks.

\begin{table}[t]
  \centering
  \caption{Known failure modes.  False negatives are cases where a broken
    algorithm passes; false positives are cases where a correct algorithm
    fails.  Most risks are mitigated by design choices or warnings.}
  \label{tab:failmodes}
  \small
  \setlength{\tabcolsep}{2.5pt}
  \begin{tabular}{@{}p{2.2cm}p{3.1cm}p{2.2cm}@{}}
    \toprule
    Source & Description & Mitigation \\
    \midrule
    \multicolumn{3}{@{}l}{\textit{False negatives}} \\[2pt]
    MaxBitDepth & Paths $>20$ bits excluded & Increase option; WARNING \\
    TimeLimit & BFS exits after 120\,s/state & WARNING printed \\
    MaxComponents & Only first $N$ components checked & Set appropriately \\
    Hidden state & Mutable globals contaminate paths & Stateless design \\
    \midrule
    \multicolumn{3}{@{}l}{\textit{False positives}} \\[2pt]
    FullSimplify incomplete & Rare; deepCheck catches most & Two-pass strategy \\
    Unlisted params & Numeric value reduces check to single point & List all in \texttt{\$checkerAbs\-tractParams} \\
    \bottomrule
  \end{tabular}
\end{table}

% ============================================================
\section{Computational Implementation}
\label{sec:mma}

\subsection{Mathematica as the host language}
\label{sec:why-mma}

The framework's central requirement is that every call to the random-number
generator can be silently intercepted and replaced by a symbolic token,
without the algorithm under test being modified or even aware of the
substitution.  Mathematica satisfies this via \texttt{Block[\{f=g,\ldots\},\,expr]},
which dynamically rebinds any named function for the duration of a lexical
scope.  Regardless of how deeply a call to \texttt{RandomReal[]} is buried
inside utility functions or conditional branches, \texttt{Block} intercepts
it transparently; the algorithm requires no cooperation.  No other mainstream
language provides this mechanism for arbitrary runtime call trees without
source modification.

The second essential feature is Mathematica's \emph{UpValue} system, which
allows comparison operators ($<$, $\leq$, \ldots) to be overloaded on the
interval-tracking token \emph{within the CAS evaluator itself}: the
expression \texttt{RandomReal[]~<~MetropolisProb[dE]} dispatches to
the symbolic token handler via UpValue without any special syntax in the
algorithm.  Together, \texttt{Block} and UpValues make the interception layer
completely invisible to the algorithm author, eliminating the most dangerous
class of false negatives: those where a random call escapes symbolic tracking
and returns a concrete floating-point number that superficially resembles a
valid symbolic weight.

For the algebraic verification step, Mathematica's \texttt{PiecewiseExpand}
reduces products of \texttt{Piecewise} functions to a flat list of cases by
computing the Cartesian product of all condition branches automatically.
\texttt{Reduce[$\mathit{cond}\wedge\beta>0$, params, Reals]}
provides a complete decision procedure for linear real arithmetic over
coupling constants.  When these are exhausted, \texttt{FullSimplify} applies
a search over hundreds of algebraic transformation rules, handling mixed
exponential-\texttt{Piecewise} expressions that arise in complex algorithms
such as VMMC with abstract coupling fields.

\medskip\noindent\textbf{Alternative implementations.}
The strongest open-source alternative is Julia combined with
Symbolics.jl~\cite{Gowda2021}.  The Cassette.jl overdubbing
framework~\cite{Cassette} provides a runtime interception mechanism close to
Mathematica's \texttt{Block}: a context can transparently redirect calls to
\texttt{rand()} anywhere in a Julia call tree to a user-defined symbolic
substitute.  A Julia implementation would remove the proprietary-software
dependency while preserving most robustness advantages.  Python with
SymPy~\cite{Meurer2017} is the weakest option: it lacks a runtime overdubbing
mechanism, SymPy's \texttt{piecewise\_fold} does not perform cross-product
expansion, and \texttt{ask} is incomplete for multivariate linear
inequalities.

\subsection{Computational cost and parallelisation}
\label{sec:parallel}

The two main computational bottlenecks are the BFS itself and the
\texttt{FullSimplify} calls.

\medskip\noindent\textbf{Parallel BFS.}
\texttt{BuildTreeAT} uses level-synchronous parallel BFS: the entire current
wave frontier is sent to \texttt{ParallelMap} in one call, partitioned into
at most $n_K$ chunks (one per available kernel).  This gives at most one
round-trip overhead per BFS wave, amortised across all states in the wave.
Single-state waves revert to sequential evaluation automatically.  On a
4-core machine, this yields a $3.3\times$ speedup for the 504-state
\texttt{vmmc\_continuous.wl} component at \texttt{NGrid=3}: sequential BFS
time 96.3\,s reduced to 28.9\,s.

\medskip\noindent\textbf{Parallel FullSimplify.}
When the fast check is inconclusive for many unique expressions, the fallback
\texttt{FullSimplify} is also parallelised.  All unique expressions requiring
\texttt{FullSimplify} are collected after syntactic deduplication and
distributed across available subkernels in one \texttt{ParallelMap} call.
The subsequent violation scan performs only fast Association lookups.  For
\texttt{vmmc\_continuous.wl} at \texttt{NGrid=3}: total check time reduced
from 29\,min to 12.5\,min on 4 kernels.

\medskip\noindent\textbf{Symmetry-aware cost savings.}
The canonical-ordering oracle and hash deduplication together achieve
near-$|G|$-fold compression.  For a two-particle component on the $3\times3$
torus (72 states, $|G|=72$ under D4+translations), the $288$ non-trivial
$(i,j)$ pairs collapse to $2$ unique expressions: 99\% of all \texttt{FullSimplify}
evaluations are eliminated.  The G-invariance check over the full 504-state
component takes $4.3$\,s on 4 kernels.  For reference, Table~\ref{tab:symtimes}
shows estimated runtimes at \texttt{NGrid=3} under different symmetry
exploitation levels (estimated from measured data).

\begin{table}[t]
  \centering
  \caption{Estimated runtimes for the \texttt{vmmc\_lattice.wl} 504-state
    component at \texttt{NGrid=3} on 4 kernels, as a function of symmetry
    exploitation.  FS = FullSimplify calls; estimates scale linearly from
    measured 12.5\,min at no-symmetry.}
  \label{tab:symtimes}
  \small
  \setlength{\tabcolsep}{3pt}
  \begin{tabular}{@{}lrrr@{}}
    \toprule
    Mode & Pairs & FS calls & Est.\ time \\
    \midrule
    No symmetry         & 126,756 & $\sim$300 & $\sim$12\,min \\
    Translations ($|G|=9$) & 14,084 & $\sim$33 & $\sim$1.4\,min \\
    D4+trans.\ ($|G|=72$) & $\approx$1,760 & $\sim$4 & $\sim$10\,s \\
    \bottomrule
  \end{tabular}
\end{table}

\subsection{User interface and algorithm files}
\label{sec:interface}

Users supply a \texttt{.wl} algorithm file defining four items: the energy
function \texttt{energy[state]}, the MC move \texttt{Algorithm[state]}, the
BitsToState decoder \texttt{BitsToState[bits]}, and a numeric inverse
temperature \texttt{numBeta} for the numerical check.  Optional definitions
control symmetry declarations, abstract parameter handling, and display
formatting.

\medskip\noindent\textbf{Abstract parameters.}
Physical parameters (e.g., \texttt{physLen}, \texttt{epsLJ}) that must be
unbound symbols during BFS are declared as string names in
\texttt{\$checkerAbstractParams}.  The checker saves their concrete values,
clears the symbols, runs BFS, then restores values for the numerical MCMC
check.  String names are required because a list of symbol references would
evaluate to their current numeric values before the checker intercepts them.

\medskip\noindent\textbf{Abstract coupling functions.}
When a coupling function \texttt{couplingJ} should have no concrete
definition during BFS (each call treated as a free real atom), a concrete
implementation \texttt{\$pairEnergy} is provided separately.  The checker
injects it via \texttt{Block[\{couplingJ\},~couplingJ~:=~\$pairEnergy]}
for numerical runs.

The checker is invoked as:
\begin{lstlisting}
wolframscript -file check.wls algo.wl \
  NGrid=N MaxComponents=N Mode=Symbolic
\end{lstlisting}
with options controlling grid size, component limits, bit depth, symmetry
exploitation, and parallelisation.  A report script (\texttt{report.wls})
provides detailed output for a single seed state, including the full decision
tree and symbolic transition matrix.

% ============================================================
\section{Results}
\label{sec:examples}

\subsection{1D Kawasaki dynamics on a ring}
\label{sec:kawasaki}

The 1D Kawasaki algorithm~\cite{Kawasaki1966} operates on a periodic ring of
$L$ sites.  Each site holds a vacancy (type 0) or one of $N_t$ particle
types.  At each step, two adjacent sites are chosen uniformly at random and
their contents swapped with Metropolis probability $\Mf{-\beta\Delta E}$,
where the energy is the nearest-neighbour sum
\begin{equation}
  E(\mathbf{s}) = \sum_{i=1}^{L} J\!\left(s_i,\, s_{i\bmod L+1},\, 1\right),
  \label{eq:kawenergy}
\end{equation}
with $J(a,b,d^2)$ a coupling function treated as a free symbol for each
distinct type pair.

For a 4-site ring with one particle of type 1 and one of type 2, the checker
discovers 12 states, which fall into two energy classes: 8 states with the
particles adjacent (energy $J_{12}$) and 4 with them separated (energy $0$).
Table~\ref{tab:kawmat} shows representative transition matrix entries.  The
detailed-balance condition for the pair
(\lsite{0}\lsite{0}\lsite{1}\lsite{2},\;\lsite{0}\lsite{1}\lsite{0}\lsite{2})
evaluates to
\begin{equation}
  \tfrac{1}{4}\Mf{\beta J_{12}}\cdot e^{-\beta J_{12}} \;=\;
  \tfrac{1}{4}\Mf{-\beta J_{12}}\cdot 1,
\end{equation}
which holds for all $J_{12}$: if $J_{12}>0$, both sides equal
$\tfrac{1}{4}e^{-\beta J_{12}}$; if $J_{12}<0$, both equal $\tfrac{1}{4}$.
All 132 non-trivial $\Delta_{ij}$ simplify to zero and the checker reports
PASS.

\begin{table}[t]
  \centering
  \caption{Representative symbolic transition matrix entries for the 4-site
    1D Kawasaki ring ($L{=}4$, particles of types 1 and 2).
    States shown as periodic lattice rows \lsite{\cdot}\lsite{\cdot}\lsite{\cdot}\lsite{\cdot}.
    All non-zero off-diagonal entries follow the same pattern.}
  \label{tab:kawmat}
  \setlength{\tabcolsep}{2pt}
  \resizebox{\columnwidth}{!}{%
  \begin{tabular}{@{}llr@{}}
    \toprule
    Transition ($i\to j$) & Description & $T(i\to j)$ \\
    \midrule
    \multicolumn{3}{@{}l}{\textit{From adjacent pair ($E = J_{12}$)}} \\[1pt]
    \lsite{0}\lsite{0}\lsite{1}\lsite{2}$\to$\lsite{0}\lsite{0}\lsite{2}\lsite{1}
      & swap types; $\Delta E=0$
      & $\tfrac{1}{4}$ \\[2pt]
    \lsite{0}\lsite{0}\lsite{1}\lsite{2}$\to$\lsite{0}\lsite{1}\lsite{0}\lsite{2}
      & separate; $\Delta E=-J_{12}$
      & $\tfrac{1}{4}\Mf{\beta J_{12}}$ \\[2pt]
    \lsite{0}\lsite{0}\lsite{1}\lsite{2}$\to$\lsite{2}\lsite{0}\lsite{1}\lsite{0}
      & separate (wrap); $\Delta E=-J_{12}$
      & $\tfrac{1}{4}\Mf{\beta J_{12}}$ \\[4pt]
    \multicolumn{3}{@{}l}{\textit{From separated pair ($E = 0$)}} \\[1pt]
    \lsite{0}\lsite{1}\lsite{0}\lsite{2}$\to$\lsite{0}\lsite{0}\lsite{1}\lsite{2}
      & bring adjacent; $\Delta E=+J_{12}$
      & $\tfrac{1}{4}\Mf{-\beta J_{12}}$ \\[2pt]
    \lsite{0}\lsite{1}\lsite{0}\lsite{2}$\to$\lsite{0}\lsite{1}\lsite{0}\lsite{2}
      & self (failed proposals)
      & $1-\Mf{-\beta J_{12}}$ \\
    \bottomrule
  \end{tabular}}
\end{table}

\subsection{Rightward-only cluster: a false negative in numerical testing}
\label{sec:cluster}

We demonstrate the diagnostic power of symbolic checking over numerical
testing using a one-dimensional cluster algorithm that always violates
detailed balance yet passes every statistical test.

The algorithm identifies all maximal connected clusters of particles on the
ring.  One cluster is chosen at random.  If the site immediately clockwise of
the cluster is vacant, the entire cluster is slid one step clockwise and
accepted with Metropolis probability $\Mf{-\beta\Delta E}$.  No leftward
slide is ever proposed.

For any transition $i\to j$ caused by a rightward slide, $T(j\to i)=0$
because the reverse leftward move is never proposed.  Thus
$\Delta_{ij} = T(i\to j)\,e^{-\beta E(i)}\neq 0$ for all $J_{ab}$ and
$\beta$.  The symbolic checker detects this immediately.

Crucially, however, the algorithm visits states in a directed cycle under
rigid cluster translation.  Because sliding a cluster rigidly preserves all
internal bonds and simply relabels the boundary bond, every state in the
cycle has the same energy.  The Boltzmann distribution over the cycle is
therefore uniform, and any numerical test comparing the empirical histogram
to the Boltzmann distribution finds perfect agreement.

Table~\ref{tab:cluster} shows checker output (Mode=Both) for several seed
states.  States where all sites are occupied form single-state trivially
balanced components.  For all multi-state components, the symbolic checker
finds violations while the numerical KL divergence is consistent with zero
and the numerical check reports PASS.

\begin{table}[t]
  \centering
  \caption{Checker output (Mode=Both) for the rightward-only 1D cluster
    algorithm.  \#St: component size.  Sym: symbolic result.  \#F: violated
    ordered pairs.  KL: Kullback--Leibler divergence from Boltzmann (100\,000
    steps).  Dash: single-state component.}
  \label{tab:cluster}
  \small
  \setlength{\tabcolsep}{3pt}
  \begin{tabular}{@{}lllrlll@{}}
    \toprule
    Bits & State & \#St & Sym. & \#F & KL & Num. \\
    \midrule
    \texttt{10}      & $\{1\}$       & 1 & PASS         & 0 & $-$  & $-$  \\
    \texttt{100}     & $\{1,0\}$     & 2 & PASS         & 0 & ${\sim}0$ & PASS \\
    \texttt{110}     & $\{1,2\}$     & 1 & PASS         & 0 & $-$  & $-$  \\
    \texttt{1001}    & $\{1,0,0\}$   & 3 & \textbf{FAIL}& 3 & ${\sim}0$ & PASS \\
    \texttt{1100}    & $\{1,2,0\}$   & 3 & \textbf{FAIL}& 3 & ${\sim}0$ & PASS \\
    \texttt{11001}   & $\{1,0,0,0\}$ & 4 & \textbf{FAIL}& 4 & ${\sim}0$ & PASS \\
    \texttt{11101}   & $\{1,2,0,0\}$ & 4 & \textbf{FAIL}& 4 & ${\sim}0$ & PASS \\
    \texttt{101001}  & $\{1,2,3,0\}$ & 4 & \textbf{FAIL}& 4 & ${\sim}0$ & PASS \\
    \bottomrule
  \end{tabular}
\end{table}

This example illustrates a broad class of latent bugs: an algorithm whose
stationary distribution is exactly Boltzmann under translation-invariant
energies, yet which violates detailed balance and would sample incorrectly
for any position-dependent energy (e.g., an external field).  No numerical
test at the symmetric energy can detect the failure; the symbolic check
catches it unconditionally, for all parameter values, without running any
Markov chain.

\subsection{2D Virtual Move Monte Carlo}
\label{sec:vmmc}

VMMC~\cite{Whitelam2007} is a cluster algorithm for particles on a 2D
periodic lattice.  A seed particle is chosen uniformly at random, a virtual
rigid displacement in a randomly chosen direction is proposed, and a cluster
is grown by recursively testing each occupied neighbour with forward and
reverse link probabilities
\begin{equation}
  p_\mathrm{fwd}(i,j) = \max\!\bigl(0,\,1-e^{\beta(\varepsilon_\mathrm{init}-\varepsilon_\mathrm{fwd})}\bigr),
\label{eq:plink}
\end{equation}
and analogously $p_\mathrm{rev}$.  A frustrated reverse link aborts the
move; otherwise the cluster translates rigidly.  This enforces detailed
balance without a Metropolis acceptance step~\cite{Whitelam2007}.

\smallskip\noindent\textbf{Correct VMMC.}
On a $2\times2$ torus with 2 distinct particles (12 states), the checker
verifies all 132 non-trivial $\Delta_{ij}$ and reports PASS.  A
representative entry is
\begin{equation}
  T\!\bigl(\{0,0,1,2\}\to\{1,2,0,0\}\bigr) = \tfrac12\bigl(1-(1-p_{12})^2\bigr),
\end{equation}
where $p_{12}=\max(0,1-e^{-\beta J_{12}})$, reflecting the two cluster
orientations through which this transition can be assembled.

\smallskip\noindent\textbf{Missing-direction variant.}
An intentionally broken variant omits the ``up'' direction.  On a $2\times2$
torus, periodic boundary conditions render up and down equivalent, so the
asymmetry is invisible and the checker correctly reports PASS.  On a
$3\times3$ torus, seeding with a two-particle state, the checker finds 144
violated pairs---the broken directional symmetry is unambiguously exposed.

\smallskip\noindent\textbf{Biased-displacement variant.}
\texttt{vmmc\_biased.wl} duplicates rightward displacements so that
$P(dx>0,dy)=2\cdot P(dx<0,dy)$.  This breaks DB (correctly caught: FAIL)
and also breaks $G$-invariance under D4 (the G-invariance check reports FAIL
for the rotate$_{90}$ generator, since a $90^\circ$ rotation maps the
over-represented rightward moves to unbiased upward moves).

\smallskip\noindent\textbf{Abstract field and coupling.}
A further variant leaves both $J(a,b,d^2)$ and an external field function
$f(x,y,L)$ completely unevaluated during the symbolic check.  Every call is
treated as an independent free real atom.  A PASS proves correctness for all
possible energy models simultaneously.  On a $2\times2$ grid with 3 distinct
particles (24 states), the checker reports PASS under these fully abstract
conditions.

\subsection{Symmetry-aware checking at NGrid=3}
\label{sec:sym-example}

\texttt{vmmc\_lattice.wl} replaces the rounded-Gaussian displacement of
\texttt{vmmc\_continuous.wl} with a uniform box
$\{(dx,dy): |dx|\leq n_\mathrm{step}, |dy|\leq n_\mathrm{step}\}\setminus\{(0,0)\}$.
This ensures $P(\mathbf{d})=P(-\mathbf{d})$ structurally and makes
all seqBernoulli weights rational constants, eliminating \texttt{erf}
functions and making the FastChecker complete (no false-positive risk from
transcendental feasibility conditions).  The algorithm declares
$\texttt{\$symmetryGroup} = \{\texttt{"translation"}, \texttt{"D4"}\}$.

At \texttt{NGrid=3} (\texttt{MaxComponents=4}, 4 kernels):
\begin{itemize}
  \item The two-particle component (72 states) produces 288 non-trivial
  $(i,j)$ pairs, which collapse via hash deduplication to \textbf{2 unique
  expressions} (99\% collapse rate).
  \item The G-invariance check over the 504-state component takes 4.3\,s.
  \item Total wall time for all checked components: $\approx 2$\,min\,55\,s.
\end{itemize}
The checker reports DB: PASS, Ergodicity: PASS (state count matches
$(S)_N$ for every component), G-invariance: PASS.

\smallskip\noindent\textbf{Ergodicity failure example.}
\texttt{kawasaki\_1d\_nonergodic.wl} demonstrates the independence of the
two checks.  Only particles of type 1 participate in swaps; particles of
types 2 and higher are frozen.  The checker reports: DB PASS (within each
invariant subspace the swap dynamics are correctly balanced), Ergodicity
FAIL (state count smaller than $(S)_N$ because frozen particles cannot
rearrange).  Detailed balance cannot detect non-ergodicity; both checks are
required for a valid sampler.

% ============================================================
\section{Intended Use Cases}
\label{sec:usecases}

\DBC{} is designed for algorithm developers and simulation practitioners in
statistical physics, soft matter, and computational chemistry.  Below we
describe the principal intended use cases and their scope.

\subsection{Verification and regression testing}
\label{sec:uc-regression}

The most immediate use is as a regression test for existing MC algorithms.
A single invocation of \texttt{check.wls} over the smallest non-trivial
system sizes (typically $2\times2$ or $3\times3$ for 2D algorithms) provides
a formal proof of correctness for all parameter values simultaneously.  This
surpasses statistical testing for two reasons: (i) a PASS is a theorem, not
a statistical outcome, and (ii) the proof covers the full parameter space,
including parameter regimes---e.g., near a phase transition, or with unusual
coupling ratios---where numerical tests would require a dedicated targeted
run.

The tool catches bugs that are invisible to numerical tests.  The rightward
cluster example (Section~\ref{sec:cluster}) is one class; others include:
asymmetric proposal distributions that happen to cancel at specific
temperatures; missing reverse moves that only matter in the presence of an
external field; energy function sign errors that are compensated by an
incorrect acceptance rule.

As a regression tool, the checker can be integrated into a continuous
integration pipeline: on any modification of the algorithm file, a fast
symbolic check (\texttt{Mode=Symbolic, MaxBitString=11111}) runs in under
two minutes and reports any newly introduced violations before code is merged.

\subsection{Algorithm development: design, certify, translate}
\label{sec:uc-design}

For new algorithm development, the natural workflow is to design and certify
in the Mathematica environment before translating to a production language.
The algorithm is implemented as a \texttt{.wl} file; the checker provides
formal verification on small systems.  Once a PASS is obtained, the algorithm
is translated to C++, Python, Julia, or Fortran for production use.

This workflow addresses a known challenge in simulation algorithm design:
detailed-balance proofs for complex cluster algorithms require careful
tracking of many forward and reverse path weights, and are susceptible to
subtle sign errors and missed edge cases.  The checker replaces this manual
book-keeping with an automated algebraic verification.

The \texttt{report.wls} script provides additional diagnostic support: for a
single seed state it prints the full BFS decision tree, the symbolic
transition matrix, and identifies all violated pairs.  This allows the
developer to understand exactly which paths contribute to each transition
probability and why the balance condition fails, guiding targeted correction.

Translation of a certified algorithm to a production language introduces a
gap: the checker proves the Mathematica implementation satisfies detailed
balance, not the C++/Python translation.  Two strategies mitigate this.
First, for simple algorithms (single-particle Metropolis, standard Kawasaki),
the subset of Mathematica used---integer arithmetic, array indexing,
conditional branches---maps mechanically to any language, and the translation
can be validated by running both on identical random seeds.  Second, for
complex algorithms, the Mathematica implementation serves as a reference: a
test suite comparing the two implementations on a comprehensive set of states
and random seeds establishes the translation's fidelity with high but not
absolute confidence.

\subsection{LLM-assisted algorithm discovery}
\label{sec:uc-llm}

A prospective use case is as the correctness oracle in a large language model
(LLM) algorithm discovery loop~\cite{FunSearch2023,AlphaEvolve2025}.  The
proposed architecture is as follows:
(1) an LLM generates a candidate MC algorithm as a Mathematica file;
(2) the checker runs symbolically and returns a binary PASS/FAIL;
(3) on FAIL, the exact symbolic expression $\Delta_{ij} = T(i\to j)\pi(i) -
T(j\to i)\pi(j)\neq 0$ for the violating state pair is fed back verbatim to
the LLM;
(4) the LLM revises the algorithm and the loop repeats.

The symbolic counterexample provides qualitatively richer feedback than
any runtime error message.  A Python traceback says ``your code crashed at
line 12''; the checker says ``the detailed-balance condition fails for states
$(3,7)$ because the expression $e^{\beta J_{12}} - e^{\beta J_{13}}$ is
nonzero''.  This is a mathematical statement about the structure of the
algorithm---identifying which part of the growth rule generated the $J_{12}$
term and which generated the $J_{13}$ term---that the LLM can reason about
algebraically to guide targeted correction.  This constitutes a qualitatively
richer feedback mode than standard execution loops in the LLM
literature~\cite{PAL2023,ReAct2023,Reflexion2023,LLMLOOP2025}, since the
oracle speaks in the same algebraic language as the problem domain rather
than simply reporting execution success or failure.

\medskip\noindent\textbf{Scope and limitations of this use case.}
This use case is prospective; no such discovery loop has yet been
implemented.  Several limitations must be noted.  Detailed balance is
\emph{sufficient} but not necessary for a Boltzmann stationary distribution:
algorithms that satisfy only global balance, such as
ECMC~\cite{Bernard2009}, would correctly fail the checker and be excluded
from consideration.  Detailed balance says nothing about ergodicity (which
must be checked separately) or about kinetic fidelity (whether the algorithm
approximates physical dynamics).  These are independent necessary conditions
for a physically useful algorithm.

More subtly, the search space for new algorithms is vast and largely
unstructured.  Unconstrained LLM generation is likely to produce many
syntactically valid but physically incoherent or non-ergodic algorithms.
The most productive approach is constrained search: give the LLM a physically
motivated algorithmic skeleton with a specific modifiable component, and ask
it to fill in that component.  Template-guided generation, combined with the
formal checker as a hard correctness constraint, mirrors the approach of
FunSearch~\cite{FunSearch2023} but with a formal mathematical oracle rather
than an empirical evaluator---a qualitative upgrade in the nature of the
correctness guarantee.

The physical kinetics target must also be operationalised.  Detailed
balance is a necessary condition for correct equilibrium sampling; it
says nothing about whether the algorithm reproduces physically relevant
dynamics (e.g., the correct relaxation timescale hierarchy or mean-squared
displacement scaling).  The physical fidelity metrics of
Section~\ref{sec:fidelity} provide a continuous scalar signal for this, but
a precise kinetic target must be specified before the loop can be used to
optimise for it.

\subsection{Physical kinetics characterisation}
\label{sec:uc-kinetics}

The physical fidelity metrics M1 and M2 (Section~\ref{sec:fidelity}) enable
quantitative comparison of the kinetic character of different MC algorithms
targeting the same equilibrium distribution.  This is useful in two
settings.

First, when choosing among multiple correct (DB-passing) algorithms for a
given model, M2 ranks their kinetic fidelity against Glauber reference
dynamics and M1 compares their relaxation timescale structure.  An algorithm
with a more negative M2 accepts a kinetic departure from Brownian dynamics in
exchange for faster decorrelation---a deliberate algorithmic acceleration.

Second, when the goal is to match the kinetics of a physical process as
closely as possible (e.g., for time-resolved comparisons with experiment or
with molecular dynamics), M2 and M1 provide a quantitative measure of how
well the MC transition matrix approximates the physical one on small tractable
systems.  These metrics are computed numerically at specific parameter values;
they do not constitute a formal proof of kinetic fidelity for all parameters.

% ============================================================
\section{Discussion}
\label{sec:discussion}

\medskip\noindent\textbf{Scope and system size.}
The exponential cost of BFS---$\mathcal{O}(N^2\cdot 2^d)$ in the number of
states $N$ and bit depth $d$---limits practical use to system sizes of roughly
$\lesssim\!10$ sites per dimension on current hardware.  This is smaller than
production simulation systems by several orders of magnitude.  However,
detailed-balance violations are a property of the algorithm, not the system
size: an algorithm that fails DB on a $2\times2$ lattice fails it on a
$200\times200$ lattice as well.  Small systems are the correct setting for
formal verification; large systems are the correct setting for production
simulation.  The workflow ``certify on small, simulate on large'' is
analogous to formal verification of small kernel functions before deploying
at scale.

\medskip\noindent\textbf{What the checker does and does not prove.}
A PASS from the symbolic checker proves, under the assumptions listed in
Section~\ref{sec:limitations}, that the Mathematica implementation of the
algorithm satisfies detailed balance for all model parameter values, on the
specific system sizes tested.  It does not prove ergodicity (checked
separately), kinetic fidelity (assessed via the M1/M2 metrics), or that a
translation of the algorithm to a production language preserves correctness.
It also does not verify correctness for all possible system sizes, only those
tested; in practice, the small system sizes that the checker can handle are
sufficient to expose all structural detailed-balance violations, since these
arise from the algorithm's move logic rather than from finite-size effects.

\medskip\noindent\textbf{Portability.}
Production MC simulations are written in C++, Fortran, Python, or Julia.
The most defensible architecture for existing codebases is a two-level
design: the production simulation runs in the native language; a reference
Mathematica implementation of the MC move is maintained alongside and
periodically tested against the production code on a comprehensive set of
states and random seeds; and the symbolic checker is applied to the reference
implementation.  For new algorithm development the natural workflow is the
reverse: design and certify in Mathematica, then translate once correctness
is established.

The clean architecture of the \texttt{.wl} algorithm files---a small,
bounded interface of four required definitions with no arbitrary I/O or
polymorphic data structures---is also compatible with a minimal domain-specific
language (DSL) approach: a single algorithm description with two
interpreters, one producing the Mathematica checker input and one producing
production C++ or Python.  Because the DSL is restricted, the two
interpreters can be verified semantically equivalent by inspection, and a
PASS constitutes a proof valid for the production code as well.  This is the
strongest possible guarantee and the most natural design for new algorithm
development.

\medskip\noindent\textbf{Comparison with model checking.}
Formal model checkers such as PRISM and Storm could in principle verify
temporal-logic properties of a Markov chain's transition matrix.  However,
they require the user to construct the chain's model by hand in a dedicated
language---effectively reproducing the case analysis of an analytic proof.
The present framework operates directly on the algorithm as executable code,
with no model-translation step, and targets the detailed-balance condition
directly as a pairwise algebraic identity.  This combination---executable
input, exact algebraic output, all parameter values simultaneously---is not
provided by any existing formal method for Markov chains.

% ============================================================
\section{Conclusions}
\label{sec:conclusion}

We have presented \DBC{}, an automated framework for the exact symbolic
verification of detailed balance in Monte Carlo algorithms.  By intercepting
all random calls and replacing them with reads from a binary tape, the
framework constructs a complete symbolic transition matrix by BFS and checks
the detailed-balance condition algebraically.  Coupling constants, field
functions, and inverse temperature are treated as free symbols throughout,
so a PASS constitutes a mathematical proof valid for all parameter values
simultaneously; a FAIL yields exact symbolic expressions identifying which
transitions violate balance.  Ergodicity is checked independently as a
complement.

The rightward-only cluster algorithm provides a striking illustration of the
method's diagnostic value: it visits every state in the correct proportion
under translation-invariant energies and passes all numerical statistical
tests, yet violates detailed balance because the reverse move is never
proposed.  No histogram-based test can detect this; the symbolic checker
catches it immediately.  The example illustrates a class of bugs latent in
algorithms for translation-invariant energies but manifest under any
position-dependent field---precisely the case not tested in symmetric systems.

Symmetry-aware checking extends the method's practical reach.  By making the
BFS produce syntactically identical expressions for $G$-orbit state pairs and
verifying $G$-invariance post-hoc using generator sweeps, we achieve near
$|G|$-fold compression of the verification cost while maintaining correctness
guarantees.  On a $3\times3$ lattice with D4+translation symmetry, this
compresses a two-particle component from 288 state pairs to 2 unique
expressions, reducing the dominant FullSimplify cost from minutes to seconds.

The primary limitation is the exponential BFS cost, restricting practical
verification to $\lesssim\!10$ sites per dimension.  This is sufficient to
certify structural correctness of the algorithm---detailed-balance violations
arise from algorithm logic, not from system size---but leaves open questions
about large-system ergodicity and kinetic properties.  Extensions under
consideration include a DSL approach for direct verified translation to
production code, and use of the checker as the formal oracle in an LLM-driven
algorithm discovery loop, where the symbolic counterexample can guide targeted
revisions in a way that runtime error feedback cannot.

% ============================================================
\appendix

\section{BitsToState: bijective state enumeration}
\label{app:bitstostate}

The reference implementation uses a bijective map between non-negative
integers and labeled-particle arrays.  An integer $\mathrm{ID}$ encodes, in
order: (i) the array length $L$, via a prefix count
$\sum_{l<L}|\mathcal{S}_l|$ where $|\mathcal{S}_l|$ is the total number of
labeled arrays of length $l$; (ii) the particle number $N$, via a secondary
prefix; (iii) the sorted positions of the $N$ particles, ranked by the
combinatorial number system; (iv) the assignment of types to positions,
ranked by the factorial number system.  Each step has an explicit inverse.

The function $\textsc{BitsToState}$ interprets a bit string as the binary
representation $\mathrm{ID}=\sum_i b_i\cdot 2^{|\mathbf{b}|-i}$ and calls
the decoder.  States are enumerated in strictly increasing order of
complexity; algorithm files may supply a \texttt{ValidStateIDs} filter that
returns null for configurations not satisfying global constraints (e.g.,
perfect-square lattice size for 2D systems), causing the checker to skip
those strings without any BFS work.

\section{Supported random call types}
\label{app:random}

\begin{table}[H]
  \centering
  \caption{Random call types supported by the checker and their
    reduction to the binary tape.}
  \label{tab:random}
  \small
  \setlength{\tabcolsep}{3pt}
  \resizebox{\columnwidth}{!}{%
  \begin{tabular}{@{}ll@{}}
    \toprule
    Call & Reduction \\
    \midrule
    \texttt{RandomReal[]} & Interval-tracking token \\
    \texttt{RandomInteger[\{lo,hi\}]} & Rejection sampling, $\lceil\log_2(h{-}l{+}1)\rceil$ bits \\
    \texttt{RandomChoice[list]} & Rejection sampling over elements \\
    \texttt{RandomChoice[w->els]} & Sequential Bernoulli decomposition \\
    \texttt{RandomVariate[Normal[$\mu$,$\sigma$]]} & Truncated Gaussian, Eq.~\eqref{eq:normalbin} \\
    \texttt{RandomPermutation[n]} & Knuth shuffle via \texttt{RandomInteger} \\
    \texttt{RandomSample[list,k]} & Partial Knuth shuffle \\
    \bottomrule
  \end{tabular}}
\end{table}

% ============================================================
\begin{thebibliography}{40}

\bibitem{Metropolis1953}
N.~Metropolis, A.~W. Rosenbluth, M.~N. Rosenbluth, A.~H. Teller, and
E.~Teller,
\textit{Equation of State Calculations by Fast Computing Machines},
J. Chem. Phys. \textbf{21}, 1087 (1953).

\bibitem{Hastings1970}
W.~K. Hastings,
\textit{Monte Carlo Sampling Methods Using Markov Chains and Their Applications},
Biometrika \textbf{57}, 97 (1970).

\bibitem{Frenkel2002}
D.~Frenkel and B.~Smit,
\textit{Understanding Molecular Simulation: From Algorithms to Applications},
2nd ed., Academic Press (2002).

\bibitem{Allen1987}
M.~P. Allen and D.~J. Tildesley,
\textit{Computer Simulation of Liquids},
Oxford University Press (1987).

\bibitem{Newman1999}
M.~E.~J. Newman and G.~T. Barkema,
\textit{Monte Carlo Methods in Statistical Physics},
Oxford University Press (1999).

\bibitem{Levin2009}
D.~A. Levin, Y.~Peres, and E.~L. Wilmer,
\textit{Markov Chains and Mixing Times},
American Mathematical Society (2009).

\bibitem{Kawasaki1966}
K.~Kawasaki,
\textit{Diffusion Constants near the Critical Point for Time-Dependent Ising
Models~I},
Phys. Rev. \textbf{145}, 224 (1966).

\bibitem{Whitelam2007}
S.~Whitelam and P.~L. Geissler,
\textit{Avoiding Unphysical Kinetic Traps in Monte Carlo Simulations of
Strongly Attractive Particles},
J. Chem. Phys. \textbf{127}, 154101 (2007).

\bibitem{Whitelam2009}
S.~Whitelam, E.~H. Feng, M.~F. Hagan, and P.~L. Geissler,
\textit{The Role of Collective Motion in Examples of Coarsening and Self-Assembly},
Soft Matter \textbf{5}, 1251 (2009).

\bibitem{SwendsenWang1987}
R.~H. Swendsen and J.-S. Wang,
\textit{Nonuniversal Critical Dynamics in Monte Carlo Simulations},
Phys. Rev. Lett. \textbf{58}, 86 (1987).

\bibitem{Wolff1989}
U.~Wolff,
\textit{Collective Monte Carlo Updating for Spin Systems},
Phys. Rev. Lett. \textbf{62}, 361 (1989).

\bibitem{LiuLuijten2004}
J.~Liu and E.~Luijten,
\textit{Rejection-Free Geometric Cluster Algorithm for Complex Fluids},
Phys. Rev. Lett. \textbf{92}, 035504 (2004).

\bibitem{Geweke2004}
J.~Geweke,
\textit{Getting it Right: Joint Distribution Tests of Posterior Simulators},
J. Am. Stat. Assoc. \textbf{99}, 799 (2004).

\bibitem{Kelly1979}
F.~P. Kelly,
\textit{Reversibility and Stochastic Networks},
Wiley (1979).

\bibitem{Norris1998}
J.~R. Norris,
\textit{Markov Chains},
Cambridge University Press (1998).

\bibitem{Manousiouthakis1999}
V.~I. Manousiouthakis and M.~W. Deem,
\textit{Strict Detailed Balance is Unnecessary in Monte Carlo Simulation},
Phys. Rev. E \textbf{59}, 2 (1999).

\bibitem{Bernard2009}
E.~P. Bernard, W.~Krauth, and D.~B. Wilson,
\textit{Event-Chain Monte Carlo Algorithms for Hard-Sphere Systems},
Phys. Rev. E \textbf{80}, 056704 (2009).

\bibitem{Glauber1963}
R.~J. Glauber,
\textit{Time-Dependent Statistics of the Ising Model},
J. Math. Phys. \textbf{4}, 294 (1963).

\bibitem{Turing1936}
A.~M. Turing,
\textit{On Computable Numbers, with an Application to the
Entscheidungsproblem},
Proc. London Math. Soc. \textbf{s2-42}, 230 (1936).

\bibitem{HornJohnson1985}
R.~A. Horn and C.~R. Johnson,
\textit{Matrix Analysis},
Cambridge University Press (1985).

\bibitem{Daws2004}
C.~Daws,
\textit{Symbolic and Parametric Model Checking of Discrete-Time Markov
  Chains},
in \textit{Proc.\ ICTAC 2004}, LNCS 3407, Springer (2005).

\bibitem{Kwiatkowska2011}
M.~Kwiatkowska, G.~Norman, and D.~Parker,
\textit{PRISM 4.0: Verification of Probabilistic Real-Time Systems},
in \textit{Proc.\ CAV 2011}, LNCS 6806, pp.~585--591, Springer (2011).

\bibitem{Dehnert2017}
C.~Dehnert, S.~Junges, J.-P.~Katoen, and M.~Volk,
\textit{A Storm is Coming: A Modern Probabilistic Model Checker},
in \textit{Proc.\ CAV 2017}, LNCS 10427, pp.~592--600, Springer (2017).

\bibitem{Gehr2016}
T.~Gehr, S.~Misailovic, and M.~Vechev,
\textit{PSI: Exact Symbolic Inference for Probabilistic Programs},
in \textit{Proc.\ PLDI 2016}, pp.~662--677, ACM (2016).

\bibitem{Meurer2017}
A.~Meurer \textit{et al.},
\textit{SymPy: Symbolic Mathematics in Python},
PeerJ Computer Science \textbf{3}, e103 (2017).

\bibitem{Bezanson2017}
J.~Bezanson, A.~Edelman, S.~Karpinski, and V.~B.~Shah,
\textit{Julia: A Fresh Approach to Numerical Computing},
SIAM Review \textbf{59}(1), 65--98 (2017).

\bibitem{Gowda2021}
S.~Gowda, Y.~Ma, A.~Cheli, M.~Gw\'{o}\'{z}d\'{z}, V.~B.~Shah,
A.~Edelman, and C.~Rackauckas,
\textit{Symbolics.jl: A Modern Computer Algebra System for a Modern
  Language},
arXiv:2105.03949 (2021).

\bibitem{Cassette}
J.~Revels,
\textit{Cassette.jl: Overdubbing Context-Dependent Execution in Julia},
\url{https://github.com/JuliaLabs/Cassette.jl} (2018).

\bibitem{Mathematica}
Wolfram Research,
\textit{Mathematica, Version 14},
Wolfram Research, Inc., Champaign, IL (2024).

\bibitem{FunSearch2023}
B.~Romera-Paredes \textit{et al.},
\textit{Mathematical Discoveries from Program Search with Large Language
  Models},
Nature \textbf{628}, 73 (2024).

\bibitem{AlphaEvolve2025}
Google DeepMind,
\textit{AlphaEvolve: A Coding Agent for Scientific and Algorithmic
  Discovery},
arXiv:2506.13131 (2025).

\bibitem{PAL2023}
L.~Gao \textit{et al.},
\textit{PAL: Program-Aided Language Models},
Proc.\ ICML, PMLR 202:10764--10799 (2023).

\bibitem{ReAct2023}
S.~Yao \textit{et al.},
\textit{ReAct: Synergizing Reasoning and Acting in Language Models},
ICLR (2023).

\bibitem{Reflexion2023}
N.~Shinn \textit{et al.},
\textit{Reflexion: Language Agents with Verbal Reinforcement Learning},
NeurIPS (2023).

\bibitem{LLMLOOP2025}
\textit{LLMLOOP: Automating Refinement of LLM-Generated Code through
  Iterative Feedback},
ICSME, arXiv:2603.23613 (2025).

\bibitem{Siepmann1992}
J.~I. Siepmann and D.~Frenkel,
\textit{Configurational Bias Monte Carlo: A New Sampling Scheme for
  Flexible Chains},
Mol. Phys. \textbf{75}, 59 (1992).

\bibitem{Hedges2015}
L.~O. Hedges,
\textit{libVMMC: Rejection-Free Monte Carlo for Hard Spheres},
\url{https://github.com/lohedges/vmmc} (2015).

\bibitem{WolframClient}
Wolfram Research,
\textit{WolframClientForPython},
\url{https://github.com/WolframResearch/WolframClientForPython} (2019).

\end{thebibliography}

\end{document}
